% GENERATED FILE. Edit canonical lessons, then run npm run generate:books.
% canonical-source-hash: fnv1a64:88d2fe5140a5d99e
% canonical-lessons: TE-C282-meyu, TE-C282-natu, TE-C282-alankarincu, TE-C282-pancu, TE-C282-sandarsincu

\chapter{To graze, To plant, To decorate, To share out, To visit}
\label{ch:te-to-graze-to-plant-to-decorate-to-share-out-to-visit}
\hlchaptermodality{282}

\begin{chapteropening}
\textbf{By the end of this chapter:} \emph{I can say to graze, to plant, to decorate, to share out and to visit.}

Complete the last lesson of chapter 282: I can say to graze, to plant, to decorate, to share out and to visit.
\end{chapteropening}

\section[\texorpdfstring{\te{మేయు} (mēyu)}{మేయు (mēyu)}]{\te{మేయు} (mēyu) — to graze}

\label{lesson:TE-C282-meyu}



\begin{quote}
Before the new one: say the Telugu for to get lost, then the Telugu for to love.
\end{quote}

\subsection*{You'll want to know: \te{మేయు}}
\textbf{\te{మేయు}} — \emph{mēyu} — \textquotedblleft{}to graze\textquotedblright{}.

\begin{grammarlens}[title={What you've built}]
One more word for this chapter.
\end{grammarlens}

\subsection*{Your turn}
\begin{itemize}
\raggedright
  \item \emph{Say it:} \emph{mēyu}
  \item \emph{Say it:} \emph{mēyu}, once more
  \item \emph{Say it:} say \emph{mēyu}
  \item \emph{Recall:} say the Telugu for to get lost, then the Telugu for to love, then say \emph{mēyu} again
\end{itemize}

\begin{tcolorbox}[breakable,colback=teal!4,colframe=teal!35!black,title={Before you move on}]
What does \te{మేయు} mean? (\textquotedblleft{}To graze\textquotedblright{}.) Say it once more.
\end{tcolorbox}
\section[\texorpdfstring{\te{నాటు} (nāṭu)}{నాటు (nāṭu)}]{\te{నాటు} (nāṭu) — to plant (a sapling)}

\label{lesson:TE-C282-natu}



\begin{quote}
Before the new one: say the Telugu for to love, then the Telugu for to graze.
\end{quote}

\subsection*{You'll want to know: \te{నాటు}}
\textbf{\te{నాటు}} — \emph{nāṭu} — \textquotedblleft{}to plant (a sapling)\textquotedblright{}.

\begin{grammarlens}[title={What you've built}]
One more word for this chapter.
\end{grammarlens}

\subsection*{Your turn}
\begin{itemize}
\raggedright
  \item \emph{Say it:} \emph{nāṭu}
  \item \emph{Say it:} \emph{nāṭu}, once more
  \item \emph{Say it:} say \emph{nāṭu}
  \item \emph{Recall:} say the Telugu for to love, then the Telugu for to graze, then say \emph{nāṭu} again
\end{itemize}

\begin{tcolorbox}[breakable,colback=teal!4,colframe=teal!35!black,title={Before you move on}]
What does \te{నాటు} mean? (\textquotedblleft{}To plant (a sapling)\textquotedblright{}.) Say it once more.
\end{tcolorbox}
\section[\texorpdfstring{\te{అలంకరించు} (alaṅkariñcu)}{అలంకరించు (alaṅkariñcu)}]{\te{అలంకరించు} (alaṅkariñcu) — to decorate}

\label{lesson:TE-C282-alankarincu}



\begin{quote}
Before the new one: say the Telugu for to graze, then the Telugu for to plant.
\end{quote}

\subsection*{You'll want to know: \te{అలంకరించు}}
\textbf{\te{అలంకరించు}} — \emph{alaṅkariñcu} — \textquotedblleft{}to decorate\textquotedblright{}.

\begin{grammarlens}[title={What you've built}]
One more word for this chapter.
\end{grammarlens}

\subsection*{Your turn}
\begin{itemize}
\raggedright
  \item \emph{Say it:} \emph{alaṅkariñcu}
  \item \emph{Say it:} \emph{alaṅkariñcu}, once more
  \item \emph{Say it:} say \emph{alaṅkariñcu}
  \item \emph{Recall:} say the Telugu for to graze, then the Telugu for to plant, then say \emph{alaṅkariñcu} again
\end{itemize}

\begin{tcolorbox}[breakable,colback=teal!4,colframe=teal!35!black,title={Before you move on}]
What does \te{అలంకరించు} mean? (\textquotedblleft{}To decorate\textquotedblright{}.) Say it once more.
\end{tcolorbox}
\section[\texorpdfstring{\te{పంచు} (pañcu)}{పంచు (pañcu)}]{\te{పంచు} (pañcu) — to share out, to distribute}

\label{lesson:TE-C282-pancu}



\begin{quote}
Before the new one: say the Telugu for to plant, then the Telugu for to decorate.
\end{quote}

\subsection*{You'll want to know: \te{పంచు}}
\textbf{\te{పంచు}} — \emph{pañcu} — \textquotedblleft{}to share out, to distribute\textquotedblright{}.

\begin{grammarlens}[title={What you've built}]
One more word for this chapter.
\end{grammarlens}

\subsection*{Your turn}
\begin{itemize}
\raggedright
  \item \emph{Say it:} \emph{pañcu}
  \item \emph{Say it:} \emph{pañcu}, once more
  \item \emph{Say it:} say \emph{pañcu}
  \item \emph{Recall:} say the Telugu for to plant, then the Telugu for to decorate, then say \emph{pañcu} again
\end{itemize}

\begin{tcolorbox}[breakable,colback=teal!4,colframe=teal!35!black,title={Before you move on}]
What does \te{పంచు} mean? (\textquotedblleft{}To share out, to distribute\textquotedblright{}.) Say it once more.
\end{tcolorbox}
\section[\texorpdfstring{\te{సందర్శించు} (sandarśiñcu)}{సందర్శించు (sandarśiñcu)}]{\te{సందర్శించు} (sandarśiñcu) — to visit}

\label{lesson:TE-C282-sandarsincu}



\begin{quote}
Before the new one: say the Telugu for to decorate, then the Telugu for to share out.
\end{quote}

\subsection*{You'll want to know: \te{సందర్శించు}}
\textbf{\te{సందర్శించు}} — \emph{sandarśiñcu} — \textquotedblleft{}to visit\textquotedblright{}.

\begin{grammarlens}[title={What you've built}]
One more word for this chapter.
\end{grammarlens}

\subsection*{Your turn}
\begin{itemize}
\raggedright
  \item \emph{Say it:} \emph{sandarśiñcu}
  \item \emph{Say it:} \emph{sandarśiñcu}, once more
  \item \emph{Say it:} say \emph{sandarśiñcu}
  \item \emph{Recall:} say the Telugu for to decorate, then the Telugu for to share out, then say \emph{sandarśiñcu} again
\end{itemize}

\begin{tcolorbox}[breakable,colback=teal!4,colframe=teal!35!black,title={Before you move on}]
What does \te{సందర్శించు} mean? (\textquotedblleft{}To visit\textquotedblright{}.) Say it once more.
\end{tcolorbox}
